\hypertarget{ej2}{}
\noindent \textbf{Ejercicio 2.} Escriba los siguientes predicados sobre números enteros en lenguaje de especificación:

\begin{enumerate}[label=\alph*)]
    \item \texttt{pred divide(x, y: $\mathbb{Z}$)} que sea verdadero si y sólo si $x$ divide a $y$ (es decir, el resto es 0).
    \item \texttt{pred sonCoprimos(x, y: $\mathbb{Z}$)} que sea verdadero si y sólo si $x$ e $y$ son coprimos.
    \item \texttt{pred mayorPrimoQueDivide(x: $\mathbb{Z}$, y: $\mathbb{Z}$)} que sea verdadero si $y$ es el mayor primo que divide a $x$.
\end{enumerate}

\vspace{0.2cm}
{\color{gray}\hrule}
\vspace{0.4cm}

\textbf{Resolución:}

\begin{enumerate}[label=\textbf{\alph*)}]

    \item \textbf{\texttt{pred divide(x: $\mathbb{Z}$, y: $\mathbb{Z}$)}} \\
    Podemos definirlo de dos maneras equivalentes. Usando la operación módulo (asegurándonos de que $x$ no sea cero para no indefinir), o usando la definición formal con el cuantificador existencial. Cualquiera de las dos es válida:
    \begin{itemize}
        \item \texttt{pred divide(x: $\mathbb{Z}$, y: $\mathbb{Z}$) \{ x $\neq$ 0 $\wedge_L$ y $\bmod$ x = 0 \}}
        \item \texttt{pred divide(x: $\mathbb{Z}$, y: $\mathbb{Z}$) \{ ($\exists$ k: $\mathbb{Z}$) (x $\times$ k = y) \}}
    \end{itemize}

    \vspace{0.4cm}
    \item \textbf{\texttt{pred sonCoprimos(x: $\mathbb{Z}$, y: $\mathbb{Z}$)}} \\
    Dos números son coprimos si el único divisor positivo que tienen en común es el 1. Acá empezamos a usar el predicado \texttt{divide} que armamos recién para modularizar:
    \begin{multline*}
    \texttt{pred sonCoprimos(x: $\mathbb{Z}$, y: $\mathbb{Z}$) \{} \\
    \texttt{($\forall$ d: $\mathbb{Z}$) (d > 0 $\wedge$ divide(d, x) $\wedge$ divide(d, y) $\rightarrow$ d = 1) \}}
    \end{multline*}

    \vspace{0.4cm}
    \item \textbf{\texttt{pred mayorPrimoQueDivide(x: $\mathbb{Z}$, y: $\mathbb{Z}$)}} \\
    Desglosamos el problema en tres condiciones que deben cumplirse a la vez:
    1. Que $y$ sea primo (usamos el predicado del Ejercicio 1).
    2. Que $y$ divida a $x$.
    3. Que cualquier otro número $p$ que también sea primo y divida a $x$, sea menor o igual a $y$ (esto garantiza que $y$ es el \textit{mayor} de todos ellos).
    
    \begin{multline*}
    \texttt{pred mayorPrimoQueDivide(x: $\mathbb{Z}$, y: $\mathbb{Z}$) \{} \\
    \texttt{esPrimo(y) $\wedge$ divide(y, x) $\wedge$} \\
    \texttt{($\forall$ p: $\mathbb{Z}$) (esPrimo(p) $\wedge$ divide(p, x) $\rightarrow$ p $\leq$ y) \}}
    \end{multline*}

\end{enumerate}

\vspace{0.5cm}
\hrule
\vspace{0.5cm}